\clearpage

\section{学位论文中的常见排版元素}

学位论文的写作离不开各类排版元素：清晰的章节层级有助于读者把握论文的逻辑骨架，公式、图与表承载着研究中的定量与可视化结果，参考文献则记录了研究的学术坐标。本章在叙述上述元素在学位论文中作用的同时，演示本模板对它们的具体处理方式，作者可对照源文件直接复制所需的代码片段。

\subsection{章节层级}

学校规范允许论文的章节最多设至四级。本模板基于ctexart文档类，依次以\verb|\section|、\verb|\subsection|、\verb|\subsubsection|与\verb|\paragraph|表示这四级标题，编号分别为“第X章”“X.Y”“X.Y.Z”与“X.Y.Z.W”，标题字体由模板统一设定。

\subsubsection{三级标题示范}

这是一段用于示范三级标题（即\verb|\subsubsection|）的正文。三级标题适用于承载具体的子主题，在大多数研究生学位论文中已足以表达论述层次。

\paragraph{四级标题示范}

这是一段用于示范四级标题（即\verb|\paragraph|）的正文。四级标题适合在论述精细到“子主题中的若干侧面”时使用。需要注意的是，四级标题不会出现在目录中，以避免目录过于臃肿。

\subsection{公式}

公式在学位论文中承担定量描述的核心职能。本模板将公式按章重置编号，编号格式为“章号-序号”，每当进入新的一章，公式编号自动归零。下式给出回归任务中常用的均方误差：
\begin{equation}
\text{MSE} = \frac{1}{n} \sum_{i=1}^{n} (y_i - \hat{y}_i)^2
\label{eq:mse}
\end{equation}
通过\verb|\label|与\verb|\eqref|配合，可在正文中引用上述公式：例如，式\eqref{eq:mse}描述了真实值与预测值之间的平均平方偏差。

\subsection{图与表}

图与表的编号规则与公式一致，均按章重置。图片建议使用矢量图或PDF文件，以保证缩放后的清晰度。
如图\ref{fig:placeholder}所示，作者只需将\verb|\includegraphics|指向自己的图片路径即可。

\begin{figure}[h]
    \centering
    \includegraphics[width=0.8\textwidth]{figure/cp1/xxx.pdf}
    \caption{图片占位示例}
    \label{fig:placeholder}
\end{figure}

复杂表格通常通过\verb|\input|从外部文件引入，以便将表格代码独立维护、避免主章节文件被表格代码淹没。表\ref{tab:xxx}的源码位于\texttt{tables/EXP.tex}，由本节通过一行\verb|\input|引入。

\input{tables/EXP}

\subsection{参考文献的引用}

中文学术写作中，参考文献的引用通常有两种形式，本模板均支持。

\textbf{行内引用}使用\verb|\cite|命令，将引用标号视为句子的语法成分，置于句中。例如：“文献\cite{cryer2011}对时间序列分析进行了系统介绍。”此时引用编号与正文同处一行，便于读者在阅读句子时即定位到引用来源。

\textbf{上标引用}使用\verb|\upcite|命令，将引用标号上提至上标位置，作为对前述论述的注解。例如：“深度学习近年来取得了显著进展\upcite{atamuradov2017prognostics}。”上标引用不打断句子的语法结构，更适合用于对一整段论述作集中标注的场景。

参考文献条目维护在项目根目录的\texttt{ref\_paper.bib}文件中，添加新条目后重新编译即可自动更新参考文献列表。\verb|ref_paper.bib|支持BibTeX的标准条目类型。除常见的\verb|@article|与\verb|@book|外，中文硕士与博士学位论文应使用\verb|@MasterThesis|与 \verb|@PhDThesis|，并填写\verb|title|、\verb|author|、\verb|school|、\verb|address|与\verb|year|等字段。本模板分别为两类条目提供了一份样例\upcite{zhu1996master,zhu1996phd}。
对于中文文献，建议显式添加\verb|language={zh}|字段。若缺少该字段，三位及以上作者的中文文献在排版时将被截断为“et al”而非“等”，与中文学术写作的惯例不符。

本模板的参考文献样式选用\verb|gbt7714-2005|而非更新的\verb|gbt7714-2015|，原因在于先前的开发者发现\verb|gbt7714-2015|对arXiv预印本条目的兼容性欠佳，遂回退至\verb|gbt7714-2005|。两者在最终排版效果上并无本质差异，作者无需为此担心格式合规性。即便如此，\verb|gbt7714-2005|对\verb|@misc|条目的支持仍然有限，直接套用arXiv官网导出的\verb|@misc|格式可能导致排版异常。若需引用arXiv文章，建议改用谷歌学术提供的\verb|@article|导出格式，本模板已在\verb|ref_paper.bib|中提供了一份样例\upcite{liu2024deepseek}，可作为模板直接复制修改。

\subsection{脚注}

学校规范要求脚注圆圈数字作为编号，以便与文后参考文献和公式编号等元素相区分。本模板在主控文件中重新定义了\verb|\thefootnote|，将\verb|\footnote|命令\footnote{这是一段示范用的脚注，用于展示圆圈数字的呈现效果。}的输出自动替换为对应的圆圈数字，作者无需关心字体的回退细节。需要注意的是，圆圈数字的字符集仅在常见符号字体中较为完整，本模板已按优先级配置了若干字体回退方案，以保证在Overleaf和本地编译时均能正常显示。受Unicode圆圈数字符号集本身的限制，本模板的圆圈数字脚注仅支持50个，超出部分将回退为普通阿拉伯数字。

\subsection{符号与缩略语}

学位论文中若大量使用专业术语的英文缩写，建议在正文之前提供“符号与缩略语说明”一节，以便读者快速查阅。本模板已通过\verb|signAndABC|环境将该节封装好，源文件位于\texttt{chapter/sign\_and\_ABC.tex}，作者可在其中直接维护一张两列的术语表，左列为缩写、右列为完整含义。表格基于\verb|longtable|实现，可在条目较多时自动跨页。

\subsection{定义、定理、引理与证明}

数学与计算机方向的学位论文常需要正式陈述定义、定理或引理。本模板在主控文件中预定义了\verb|definition|、\verb|theorem|、\verb|lemma|、\verb|corollary|与\verb|proposition|五个环境，并通过\verb|amsthm|宏包提供\verb|proof|证明环境。下面以半群为例分别给出示范：

\begin{definition}
    设$X$为非空集合，若$X$上存在满足结合律的二元运算$\cdot:X\times X\to X$，则称$(X,\cdot)$为一个半群。
\end{definition}

\begin{proposition}
    半群$(X,\cdot)$是幺半群当且仅当其含有单位元。
\end{proposition}

\begin{lemma}
    \label{lem:periodic}
    设$(X,\cdot)$为有限半群，则对任意$a\in X$，存在正整数$m,n$使得$a^m=a^{m+n}$。
\end{lemma}

\begin{theorem}
    \label{thm:idempotent}
    任一有限半群中均存在幂等元。
\end{theorem}

\begin{proof}
    由引理\ref{lem:periodic}，对任意$a\in X$存在正整数$m,n$使得$a^m=a^{m+n}$。取$k$满足$kn\geq m$，则$a^{kn}\cdot a^{kn}=a^{2kn}=a^{kn}$，即$a^{kn}$为$X$中的幂等元。
\end{proof}

\begin{corollary}
    在有限群中，唯一的幂等元为单位元。
\end{corollary}

在编号上，\verb|lemma|与\verb|corollary|和\verb|theorem|共享同一计数器，三者按出现顺序连续编号；\verb|definition|与\verb|proposition|则各自独立计数。\verb|proof|环境会在结尾自动追加证明终止符$\square$，作者无需手工标注。
